\documentclass[12pt, a4paper]{article}

% --- Pacotes Básicos ---
\usepackage[utf8]{inputenc}
\usepackage[T1]{fontenc}
\usepackage[main=brazilian, english]{babel}
\usepackage{indentfirst}
\usepackage{geometry}
\geometry{left=3cm, top=3cm, right=2cm, bottom=2cm}

% --- Informações do Documento ---
\title{Resenha: O Mundo Global visto do lado de cá}
\author{Erickson Giesel Müller}
\date{\today}

\begin{document}
\maketitle
\thispagestyle{empty} % primeira página não tem número.

	O documentário "Encontro com Milton Santos: o Mundo Global visto do lado de cá" é narrado a partir de uma entrevista com o sociólogo brasileiro, que possui grande renome na área da geografia, cujo tema principal é o efeito colateral que a globalização teve nos países de terceiro mundo. O autor do filme separa o mundo em 3 tipos: o mundo como ele é, o mundo como fábula e o mundo como ele pode ser, assim, 3 tipos diferentes de globalização.

	Antes de se aprofundar nessa tese, o filme exibe as consequênciasa do consenso de Washington na América Latina, resultando em manifestações populares que se opunham à vontade do primeiro mundo de transformar a América Latina num modelo neoliberal. O então definido "século das revoluções" se encerrou com inúmeras manifestações iniciadas para serem resolvidas no próximo século. A população dos países emergentes estavam perdendo seu senso de identidade, a globalização trouxe o consumismo ao demográfico mais pobre antes de resolver problemas como a fome e a desigualdade social.

	Pode-se dizer que a globalização agravou a desigualdade social no lado de cá, com o fim da guerra fria, e a abertura econômica da China, o ser humano deixou de ser o centro do mundo e este passou a ser o dinheiro. Por conseguinte, as políticas econômicas e sociais passaram a dar um peso maior ao dinheiro em suas políticas, deixando o humano de lado. Em certo momento do documentário, podemos ver o discurso do até então representante do Banco Mundial defendendo a privatização de água, para que esta deixe de ser considerada um direito da humanidade.

	Com a desestruturalização social que a globalização segue fazendo na organização popular, novas fontes de expressões nascem, sobretudo na periferia, que buscam utilizar das ferramentas do capitalismo para registrar sua posição no mundo. O documentário traz excertos de diversas pessoas que trabalham no audiovisual, produzindo cinema de guerrilha com suas câmeras e equipamento de baixo custo. Essas expressões são uma resposta agressiva às políticas globais para o terceiro mundo, a população emergente não quer consumir aquilo que a classe dominante quer que ela consuma, mas sim produzir a sua própria cultura para que seja apreciada pela própria população.

	Milton Santos vê essas manifestações com bastante otimismo, segundo ele, essa conquista seria possível de quebrar a opressão que o hemisfério norte pratica. Contudo, vinte anos depois do filme ser lançado, podemos perceber certa ingenuidade na visão do geógrafo. A população terceiromundista tem sim completo acesso às ferramentas do capitalismo, mas esse acesso deliberado nunca prometeu que fosse permitir que as pessoas do século XXI pudessem usá-lo para expressar seu lugar no mundo. Em vez disso, podemos perceber que a periferia do mundo está cada vez mais consumista e a desigualdade social continua crescendo exponencialmente.

\end{document}
